\section{Introduction}

Camera relocalization estimates the six-degree-of-freedom pose of a query image in a known scene. Scene coordinate regression (SCR) methods solve this problem by predicting dense 3D scene coordinates and then recovering camera pose with a geometric solver. Recent systems such as ACE and GLACE have made this pipeline fast and accurate by improving coordinate encoders and global-local scene representations. However, these methods still face a central representation question: how should a relocalization model store and reuse scene-specific memory without tying that memory to a single backbone or relying on global features that may be unreliable in some scenes?

We address this question with \method{} (\shortmethod{}), a plug-in module that compresses scene memory into compact latent tokens and injects the resulting representation into relocalization backbones. Instead of treating scene memory as a dense map that must be consumed directly, \shortmethod{} learns a compressed interface between a memory source and a coordinate-regression model. The memory source can be geometry-aware features from MapAnything, feature-space memory extracted from a frozen ACE-FCN encoder, or global-local representations used by GLACE-style regressors. This design makes \shortmethod{} a representation enhancement rather than a replacement for SCR backbones.

Our experiments follow two deployment patterns. In the one-level two-stage setting, DINOv2 query features are conditioned on MapAnything-derived memory and trained through memory alignment followed by reprojection refinement. This setting gives the strongest current Indoor6 evidence in our codebase. In the two-level two-stage setting, ACE-FCN memory first trains a local latent scene representation, and a second stage connects this local stack to a GLACE-style global-local regressor. This setting tests whether \shortmethod{} remains useful beyond the DINOv2/MapAnything configuration.

A key observation is that global conditioning is not uniformly reliable. It can improve compatible scenes, such as Wayspots Bears, while causing negative transfer in scenes such as SquareBench. We therefore define a visibility-guided configuration selection rule, denoted \gvcs{}, that estimates memory visibility and selects global, local, or hierarchical conditioning accordingly. In the current implementation, \gvcs{} is a deterministic routing rule based on front-facing memory visibility, not a learned policy. This distinction is important because it limits the claim to reliability-guided configuration selection rather than universal automatic global-feature selection.

This paper makes three contributions. First, we introduce \shortmethod{} as a compact latent scene-memory interface for plug-in relocalization conditioning. Second, we instantiate the same principle in both one-level and two-level two-stage training pipelines, covering DINOv2 with MapAnything memory and ACE-FCN memory with GLACE-style global-local regression. Third, we report evidence and limitations across Indoor6 and Wayspots, including positive gains and negative-transfer cases, and outline the remaining experiments needed to complete the AAAI submission.
